\subsection{\texorpdfstring{\(\overline{\mathbf{R}}\) 中的可和族与可乘族}{SUMMABLE FAMILIES AND MULTIPLIABLE FAMILIES IN \textbackslash overline\{\textbackslash mathbf\{R\}\}}}

在区间 \([0, +\infty]\) （\(\overline{\mathbf{R}}\) 中的）上，加法是一个结合的、交换的合成法则（§ 4，第 3 节）；因此在该区间中数的可和族的概念已有定义（第三章，§ 5，第 1 节，注 3）。

命题 2. 族 \((x_{\mathfrak{t}})\) （由实数 \(\geqslant 0\) 组成）的每一个在 \(\overline{\mathbf{R}}\) 中可和。因为映射 \(\mathbf{H} \to s_{\mathbf{H}}\) （由有向集 \(\mathfrak{F}(\mathbf{I})\) 到 \(\overline{\mathbf{R}}\) 中）是递增的；故（§ 5，第 2 节，定理 2）有极限。

同样的推理表明，实数 \(\leqslant 0\) 的每个族在 \(\overline{\mathbf{R}}\) 中可和。

类似地，乘法在区间 \([0, 1]\) 与 \([1, +\infty]\) （\(\overline{R}\) 中的）的每一个上都是一个结合的、交换的合成法则，因此可乘族的概念在每个这样的区间上都有定义。

命题 3. 每个族 \((1 + u_{t})\) {[}相应地，\((1 - u_{t})]\) ——由数 \(\geqslant 1\) （相应地，\(\geqslant 0\) 且 \(\leqslant 1\) ）组成——在 \(\overline{\mathbf{R}}\) 中可乘。

与命题 2 的证明相同。

推论. 积 \(\prod_{i} (1 + u_i) [\text{resp. } \prod_{i} (1 - u_i))\) （由数 \(\geqslant 1\) （相应地，\(>0\) 且 \(\leqslant 1\) ）组成）等于 \(+\infty\) （相应地，0）当且仅当 \(\sum_{i} u_i = +\infty\) 。因为若 \(\sum_{i} u_i\) 有限，则 \(\prod_{i} (1 + u_i)\) 与 \(\prod_{i} (1 - u_i)\) 属于 \(\mathbf{R}^*\) ，反之亦然（第 4 节，定理 4）。

注. 结合性定理（第三章，§ 5，第 3 节，定理 2）当把 G 换成 \(\overline{\mathbf{R}}\) 并假设诸 \(x_{i}\) \(\geqslant 0\) 时仍然成立。若 \(\sum_{i \in I} x_{i}\) 有限，这是显然的。反之设 \(\sum_{i \in I} x_{i} = +\infty\) 。则对每个有限的 \(a > 0\) ，存在 I 的有限子集 H 使得 \(\sum_{i \in H} x_{i} \geqslant a\) 。设 K 是 L 的有限子集且 \(H \subset \bigcup_{\lambda \in K} I_{\lambda}\) ；那么由于 \(s_{\lambda} \geqslant \sum_{i \in I_{\lambda} \cap H} x_{i}\) 对一切 \(\lambda \in K\) ，我们有

\[
\sum_ {\lambda \in \mathbf {K}} s _ {\lambda} \geqslant \sum_ {i \in \mathbf {H}} x _ {i} \geqslant a,
\]

这表明 \(\sum_{\lambda \in L} s_{\lambda} = +\infty\) 。我们把对 \([0, 1]\) 或 \([1, +\infty]\) 中的数的可乘族表述类似命题的任务留给读者。

